\documentclass[12pt, a4paper]{article}

% ── Pacotes ─────────────────────────────────────────────────────
\usepackage[utf8]{inputenc}
\usepackage[T1]{fontenc}
\usepackage[brazil]{babel}
\usepackage[top=2.5cm, bottom=2.5cm, left=2.5cm, right=2.5cm]{geometry}
\usepackage{setspace}
\usepackage{indentfirst}
\usepackage{titlesec}
\usepackage{enumitem}
\usepackage{hyperref}
\usepackage{parskip}
\usepackage{times}
\usepackage{booktabs}
\usepackage{array}

\onehalfspacing
\setlength{\parindent}{0pt}
\setlength{\parskip}{6pt}

\setlist[itemize,1]{label=\textbullet, leftmargin=1.5em}
\setlist[itemize,2]{label=\textbullet, leftmargin=1.5em}

\titleformat{\section}{\normalfont\bfseries\large}{\thesection.}{0.5em}{}
\titleformat{\subsection}{\normalfont\bfseries}{\thesubsection.}{0.5em}{}
\titlespacing{\section}{0pt}{12pt}{4pt}
\titlespacing{\subsection}{0pt}{8pt}{2pt}

\hypersetup{colorlinks=true, linkcolor=black, citecolor=black, urlcolor=blue}

% ================================================================
\begin{document}
% ================================================================

% ── Cabeçalho ────────────────────────────────────────────────────
\begin{center}
{\large\bfseries AA02: Identificação \& Proposta de Metodologia}\\[0.3em]
\textit{Disciplina: Metodologia de Pesquisa}\\
\textit{Profa: Glívia Angélica Rodrigues Barbosa}\\
\textit{Departamento de Computação (DECOM)}\\
\textit{Centro Federal de Educação Tecnológica de Minas Gerais (CEFET-MG)}
\end{center}

\vspace{0.8em}
\underline{\textit{\textbf{Nome:} Morais, Ester}}

\vspace{1em}

% ================================================================
\section{Caracterização da Metodologia dos Trabalhos Relacionados}
% ================================================================

% ── Trabalho 1 ──────────────────────────────────────────────────
\subsection{Trabalho: PentestGPT: An LLM-Empowered Automatic Penetration Testing Framework}

\begin{itemize}

  \item \textbf{Referência:}
  DENG, Gelei; LIU, Yi; MAYORAL-VILCHES, Victor; LIU, Peng; LI, Yuekang; XU, Yuan; ZHANG, Tianwei; LIU, Yang; PINZGER, Martin; RASS, Stefan.
  PentestGPT: An LLM-Empowered Automatic Penetration Testing Framework.
  In: \textit{Proceedings of the USENIX Security Symposium}, 2024.
  Disponível em: \url{https://arxiv.org/abs/2308.06782}.

  \item \textbf{Caracterização da metodologia:}
  \begin{itemize}

    \item[(a)] \textbf{Natureza:} Aplicada -- o trabalho parte de um problema real da indústria (custo e escassez de especialistas em \textit{pentesting}) e produz um artefato tecnológico operacional (o \textit{framework} PentestGPT) como resposta direta a essa demanda (DENG et al., 2024).

    \item[(b)] \textbf{Abordagem:} Quali-quantitativa. A avaliação utiliza métricas quantitativas (percentual de tarefas de \textit{pentesting} completadas na plataforma HackTheBox -- HTB) e, complementarmente, análise qualitativa dos padrões de erro do modelo, identificando categorias de falha recorrentes.

    \item[(c)] \textbf{Etapas -- 5 (cinco):}
    (1)~Revisão da literatura sobre \textit{pentesting} manual e uso de LLMs em segurança;
    (2)~\textit{Design} da arquitetura em três módulos (\textit{Reasoning}, \textit{Generation} e \textit{Parsing});
    (3)~Implementação do PentestGPT integrado ao GPT-4;
    (4)~Definição do \textit{benchmark} de avaliação a partir de desafios da plataforma HTB em dificuldades variadas;
    (5)~Execução do experimento e análise comparativa frente a \textit{baselines} (LLM sem estruturação de \textit{pentesting}).

    \item[(d)] \textbf{Procedimento(s) adotado(s):}
    Estudo experimental controlado com \textit{benchmark} padronizado (HTB). Os autores executam o PentestGPT e os \textit{baselines} sobre o mesmo conjunto de desafios e comparam as taxas de completude de tarefas. Análise qualitativa complementar é realizada para categorizar os modos de falha.

    \item[(e)] \textbf{Métodos e/ou técnicas utilizadas:}
    \textit{Prompting} com GPT-4 (poucas amostras); decomposição de tarefas por módulos especializados; avaliação em CTF (\textit{Capture the Flag}) como \textit{proxy} de cenários reais de \textit{pentesting}; comparação entre grupo experimental (PentestGPT) e grupo controle (LLM sem estrutura).

    \item[(f)] \textbf{Tipo(s) de resultado(s) gerado(s):}
    (i)~\textit{Framework} operacional (PentestGPT);
    (ii)~métricas quantitativas de desempenho (taxa de completude de tarefas HTB por nível de dificuldade);
    (iii)~taxonomia qualitativa dos padrões de erro do modelo;
    (iv)~\textit{benchmark} de avaliação de LLMs em \textit{pentesting} (disponibilizado publicamente).

    \item[(g)] \textbf{Limitações da iniciativa:}
    O \textit{benchmark} restringe-se a desafios HTB no estilo CTF, que diferem substancialmente de ambientes de produção reais -- especialmente SPAs com lógica de negócio.
    O trabalho não aborda vulnerabilidades semânticas (lógica de negócio, IDOR contextual) nem o desafio de navegar interfaces geradas dinamicamente por JavaScript.
    A dependência do GPT-4 (modelo proprietário) limita a reprodutibilidade e o controle experimental.
    O tamanho da amostra de desafios não é suficiente para inferência estatística formal.

  \end{itemize}
\end{itemize}

% ── Trabalho 2 ──────────────────────────────────────────────────
\subsection{Trabalho: ReAct: Synergizing Reasoning and Acting in Language Models}

\begin{itemize}

  \item \textbf{Referência:}
  YAO, Shunyu; ZHAO, Jeffrey; YU, Dian; DU, Nan; SHAFRAN, Izhak; NARASIMHAN, Karthik; CAO, Yuan.
  ReAct: Synergizing Reasoning and Acting in Language Models.
  In: \textit{Proceedings of the International Conference on Learning Representations (ICLR)}, 2023.
  Disponível em: \url{https://arxiv.org/abs/2210.03629}.

  \begin{itemize}

    \item[(a)] \textbf{Natureza:} Básica/fundamental -- o trabalho não parte de uma demanda de aplicação imediata, mas de uma questão teórica sobre a capacidade dos LLMs de integrar raciocínio e ação. O resultado (o padrão ReAct) é uma contribuição ao campo de NLP que serve de fundação para trabalhos aplicados subsequentes (YAO et al., 2023).

    \item[(b)] \textbf{Abordagem:} Quantitativa. Os experimentos utilizam \textit{benchmarks} padronizados com métricas objetivas (\textit{Exact Match}, F1-\textit{Score}, taxa de sucesso em tarefas de tomada de decisão), permitindo comparação direta entre o método proposto e os \textit{baselines}.

    \item[(c)] \textbf{Etapas -- 4 (quatro):}
    (1)~Identificação do problema: ausência de integração entre traços de raciocínio e ações externas nos LLMs existentes;
    (2)~Formulação do padrão ReAct como estratégia de \textit{prompting} que intercala \textit{thoughts} (raciocínio) e \textit{actions} (ações sobre ambiente);
    (3)~Implementação e avaliação em quatro \textit{benchmarks} de naturezas distintas (HotpotQA, FEVER, ALFWorld, WebShop);
    (4)~Análise comparativa e estudo de ablação (\textit{ReAct vs. CoT-only vs. Act-only}).

    \item[(d)] \textbf{Procedimento(s) adotado(s):}
    \textit{Few-shot prompting} com o modelo PaLM (540B parâmetros) sobre \textit{benchmarks} publicamente disponíveis.
    Para cada tarefa, os autores constroem exemplos de demonstração (\textit{few-shot}) que ilustram o padrão Thought--Action--Observation.
    O estudo de ablação isola a contribuição de cada componente do padrão (raciocínio puro, ação pura, e ambos combinados).

    \item[(e)] \textbf{Métodos e/ou técnicas utilizadas:}
    \textit{Few-shot prompting}; \textit{Chain-of-Thought} como \textit{baseline} de raciocínio; interação com ambientes simulados (ALFWorld, WebShop) via APIs; métricas de \textit{Exact Match} e F1 para tarefas de QA; taxa de sucesso para tarefas de navegação.

    \item[(f)] \textbf{Tipo(s) de resultado(s) gerado(s):}
    (i)~Padrão de \textit{prompt} ReAct (Thought--Action--Observation), publicado e replicável;
    (ii)~resultados quantitativos em quatro \textit{benchmarks}, com ganhos de até 34\% sobre os \textit{baselines} em HotpotQA;
    (iii)~análise qualitativa de exemplos ilustrando as vantagens de interpretabilidade do padrão.

    \item[(g)] \textbf{Limitações da iniciativa:}
    Os experimentos utilizam o PaLM (540B parâmetros), modelo que não estava disponível publicamente à época da publicação, o que restringe a reprodutibilidade independente.
    O trabalho avalia o padrão exclusivamente em tarefas de NLP e navegação simulada, sem validação em domínios de segurança de \textit{software} ou automação de navegador real.
    O custo computacional dos experimentos é elevado, tornando a replicação acessível apenas a laboratórios com recursos significativos.

  \end{itemize}
\end{itemize}

% ================================================================
\section{Proposta de Metodologia do Projeto}
% ================================================================

\textbf{Projeto:} Agentic DAST -- \textit{Framework} de Detecção Automatizada de Vulnerabilidades em \textit{Single Page Applications} via Agentes LLM

\vspace{0.5em}

\begin{itemize}

  \item \textbf{Questão de pesquisa a ser explorada:}
  Em que medida um agente baseado em LLM, equipado com automação de navegador e orientado pelo padrão de raciocínio ReAct, é capaz de detectar vulnerabilidades em \textit{Single Page Applications} com maior eficácia (Precisão, Revocação e F1-\textit{Score}) do que ferramentas DAST convencionais?

  \item \textbf{Natureza da iniciativa:} Aplicada -- a pesquisa parte de uma limitação concreta de ferramentas DAST em lidar com SPAs e produz um artefato tecnológico (o \textit{framework} Agentic DAST) cuja avaliação empírica valida as afirmações de desempenho (WAZLAWICK, 2014).

  \item \textbf{Abordagem metodológica:} Quantitativa -- os resultados são medidos por Precisão, Revocação e F1-\textit{Score}, calculados a partir de uma matriz \textit{Ground Truth} construída sobre o catálogo oficial de desafios do OWASP Juice Shop e comparados entre três cenários experimentais por meio de teste estatístico não-paramétrico.

  \item \textbf{Fases/Etapas da metodologia -- 6 (seis):}

  \vspace{0.5em}

  \textbf{Fase 1 -- Revisão da Literatura}
  \begin{itemize}
    \item \textit{Objetivo (o quê e por quê?):} Mapear o estado da arte em DAST, segurança de SPAs, LLMs aplicados à segurança e agentes de IA, identificando lacunas que a proposta deve preencher e fundamentando as escolhas de design.
    \item \textit{Como será conduzida:} Busca estruturada nas bases ACM Digital Library, IEEE Xplore e USENIX com \textit{strings} combinando termos como \textit{``DAST''}, \textit{``LLM security''}, \textit{``SPA vulnerability''} e \textit{``agentic AI''}. Leitura crítica e síntese em tabela comparativa de trabalhos relacionados.
    \item \textit{Tipo de resultado que gera:} Tabela de trabalhos relacionados com dimensões de comparação (automação, semântica, suporte a SPA, categorias de vulnerabilidade); identificação formal das lacunas que motivam a proposta.
  \end{itemize}

  \vspace{0.4em}
  \textbf{Fase 2 -- Design e Especificação do Framework}
  \begin{itemize}
    \item \textit{Objetivo:} Definir a arquitetura do Agentic DAST em módulos coesos e responsabilidades bem delimitadas, antes da implementação, de modo a garantir que o artefato seja extensível e avaliável.
    \item \textit{Como será conduzida:} Especificação dos quatro módulos -- (1)~\textbf{Percepção} (\textit{DOMExtractor} via Puppeteer/CDP), (2)~\textbf{Cognição} (\textit{ReActAgent} com Anthropic Claude), (3)~\textbf{Execução} (\textit{ActionDispatcher} + Casos de Abuso) e (4)~\textbf{Auditoria} (\textit{ZapClient}) -- com diagramas de componentes e definição das interfaces de dados (\textit{AppMap}, \textit{Finding}, \textit{AgentResult}).
    \item \textit{Tipo de resultado que gera:} Documento de arquitetura com diagrama de fluxo e especificação formal dos oito Casos de Abuso (categorias: Lógica de Negócio, IDOR, XSS, JWT, Dados Sensíveis, SQL Injection).
  \end{itemize}

  \vspace{0.4em}
  \textbf{Fase 3 -- Implementação}
  \begin{itemize}
    \item \textit{Objetivo:} Produzir um artefato funcional que realize o design especificado na fase anterior, permitindo sua execução reproduzível em ambiente controlado.
    \item \textit{Como será conduzida:} Desenvolvimento em TypeScript/Node.js, com controle de versão no GitHub. Implementação incremental por módulo: (i)~DOMExtractor com extração de rotas e formulários Angular; (ii)~ReActAgent com loop Thought--Action--Observation sobre a API da Anthropic; (iii)~oito funções de Abuso pré-definidas; (iv)~ZapClient para coleta de alertas passivos do OWASP ZAP.
    \item \textit{Tipo de resultado que gera:} Código-fonte funcional, versionado e publicamente disponível, com scripts de execução para os três cenários experimentais (\texttt{npm run scenario:a/b/c}).
  \end{itemize}

  \vspace{0.4em}
  \textbf{Fase 4 -- Construção do Ground Truth e Protocolo Experimental}
  \begin{itemize}
    \item \textit{Objetivo:} Definir a verdade de referência que permitirá calcular Precisão, Revocação e F1-\textit{Score} de forma objetiva e reproduzível, evitando subjetividade na avaliação.
    \item \textit{Como será conduzida:} Mapeamento manual dos desafios do OWASP Juice Shop (catálogo oficial) (KIMMINICH, 2024) para as seis categorias de vulnerabilidade cobertas pelo \textit{framework}. Para cada vulnerabilidade, registra-se o identificador do desafio, a categoria, o \textit{endpoint} afetado e o critério de detecção (resposta HTTP esperada, payload, código de status). Definição das hipóteses: H1~(Agentic DAST $>$ ZAP spider), H2~(Agentic DAST $>$ ZAP + E2E) e H3~(distribuição de categorias detectadas difere entre cenários).
    \item \textit{Tipo de resultado que gera:} Matriz \textit{Ground Truth} com $k$ vulnerabilidades; protocolo de execução com 3 rodadas por cenário ($n = 3$); critérios formais de TP, FP e FN para cada \textit{finding}.
  \end{itemize}

  \vspace{0.4em}
  \textbf{Fase 5 -- Execução do Experimento Comparativo}
  \begin{itemize}
    \item \textit{Objetivo:} Coletar dados empíricos comparando os três cenários para responder à questão de pesquisa com evidências mensuráveis.
    \item \textit{Como será conduzida:} Execução de três cenários em ambiente Docker controlado (Juice Shop + ZAP em contêineres isolados): \textbf{Cenário A} -- OWASP ZAP com \textit{spider} tradicional + \textit{scan} ativo (baseline automático); \textbf{Cenário B} -- OWASP ZAP guiado por script de navegação E2E com Puppeteer (baseline guiado por humano); \textbf{Cenário C} -- Agentic DAST autônomo com loop ReAct (YAO et al., 2023) + Casos de Abuso (experimental). As vulnerabilidades alvo são derivadas do OWASP Top 10 (OWASP FOUNDATION, 2021). Cada cenário é executado $n = 3$ vezes. Os \textit{findings} de cada execução são classificados manualmente como TP, FP ou FN com base na matriz \textit{Ground Truth}.
    \item \textit{Tipo de resultado que gera:} Relatórios JSON padronizados por execução; tabelas de Precisão, Revocação e F1-\textit{Score} por cenário e por categoria de vulnerabilidade; tempo de execução de cada cenário.
  \end{itemize}

  \vspace{0.4em}
  \textbf{Fase 6 -- Análise Estatística e Redação}
  \begin{itemize}
    \item \textit{Objetivo:} Determinar se as diferenças de F1-\textit{Score} entre os cenários são estatisticamente significativas e documentar os resultados na forma de artigo científico.
    \item \textit{Como será conduzida:} Teste de Wilcoxon \textit{signed-rank} (não-paramétrico, adequado a amostras pequenas com $n = 3$) com $\alpha = 0{,}05$ para cada par de cenários (A$\times$C, B$\times$C). Redação do artigo no formato de conferência brasileira (SBSEG ou SBES), com seções de Introdução, Trabalhos Relacionados, \textit{Framework}, Avaliação e Conclusão.
    \item \textit{Tipo de resultado que gera:} $p$-valores e decisão sobre H1, H2 e H3; artigo científico completo submetível a conferência da área.
  \end{itemize}

  \item \textbf{Limitações:}
  \begin{itemize}
    \item \textbf{Validade externa restrita:} a avaliação é realizada em um único \textit{testbed} (OWASP Juice Shop), o que limita a generalização dos resultados para outras aplicações e domínios.
    \item \textbf{Amostra pequena:} $n = 3$ rodadas por cenário é o mínimo viável para o teste de Wilcoxon; a amostra não permite inferências com alta potência estatística.
    \item \textbf{Não determinismo do LLM:} o agente produz saídas diferentes entre execuções para a mesma entrada, o que introduz variabilidade não controlada no Cenário C.
    \item \textbf{Dependência de API proprietária:} o Cenário C utiliza a API da Anthropic (Claude), cuja disponibilidade, custo e versão de modelo podem afetar a reprodutibilidade do experimento.
    \item \textbf{Cobertura parcial do OWASP Top 10:} o \textit{framework} cobre seis categorias; vulnerabilidades como CSRF, \textit{Insecure Deserialization} e falhas de configuração de segurança permanecem fora do escopo.
  \end{itemize}

  \item \textbf{Justificativa da metodologia proposta:}
  A escolha do experimento controlado com três cenários comparativos é a mais adequada para isolar o efeito do componente LLM sobre a eficácia da detecção.
  O Cenário A representa o limite inferior (DAST sem qualquer guia semântico), o Cenário B isola o efeito da navegação humana (sem LLM), e o Cenário C combina navegação autônoma e raciocínio semântico -- permitindo atribuir causalmente os ganhos de F1 à contribuição do agente.
  O uso do OWASP Juice Shop como \textit{testbed} garante reprodutibilidade e comparabilidade com outros trabalhos da literatura que utilizam o mesmo alvo.
  O teste de Wilcoxon é o método estatístico mais adequado para o tamanho de amostra ($n = 3$) e para a natureza ordinal das métricas, dispensando a premissa de normalidade que o teste $t$ exigiria.
  A abordagem de \textit{Design Science} (Hevner et al., 2004) subjacente à pesquisa tecnológica garante que o artefato seja simultaneamente útil (avaliado empiricamente) e contribuição ao conhecimento científico (publicado e replicável).

\end{itemize}

% ================================================================
\section{Conclusões sobre a Atividade Avaliativa}
% ================================================================

A elaboração desta atividade foi fundamental para consolidar dois aspectos da proposta que, até então, permaneciam implícitos no trabalho de desenvolvimento.

O primeiro diz respeito à \textbf{granularidade das fases metodológicas}.
Ao ser exigida a descrição explícita de cada etapa -- com objetivo, procedimento e tipo de resultado -- ficou evidente que algumas fases do projeto estavam sendo conduzidas de forma intuitiva, sem critério formal de conclusão.
Em particular, a Fase 4 (construção do \textit{Ground Truth}) revelou-se crítica: sem uma matriz de referência rigorosa, qualquer métrica de Precisão e Revocação torna-se subjetiva e não comparável.
A atividade forçou a formalização desse protocolo antes que o experimento fosse executado.

O segundo aspecto refere-se à \textbf{análise comparativa das metodologias dos trabalhos relacionados}.
A caracterização sistemática do PentestGPT e do ReAct evidenciou que ambos adotam avaliações experimentais -- porém com \textit{benchmarks} e domínios distintos do proposto neste trabalho.
Esse exercício confirmou que a lacuna não é apenas temática (SPAs vs.\ infraestrutura), mas também metodológica: nenhum dos trabalhos revisados realiza experimentos com grupos de controle que isolem o efeito do componente LLM sobre a detecção de vulnerabilidades semânticas.
O desenho dos três cenários experimentais (A, B e C) foi, portanto, diretamente informado por essa análise.

% ================================================================
\section{Referências}
% ================================================================

\begin{itemize}

  \item DENG, Gelei et al. PentestGPT: An LLM-Empowered Automatic Penetration Testing Framework. In: \textit{Proceedings of the USENIX Security Symposium}, 2024. Disponível em: \url{https://arxiv.org/abs/2308.06782}.

  \item HEVNER, Alan R.; MARCH, Salvatore T.; PARK, Jinsoo; RAM, Sudha. Design Science in Information Systems Research. \textit{MIS Quarterly}, v.~28, n.~1, p.~75--105, 2004.

  \item KIMMINICH, Bjoern. OWASP Juice Shop: Probably the most modern and sophisticated insecure web application. GitHub, 2024. Disponível em: \url{https://github.com/juice-shop/juice-shop}.

  \item OWASP FOUNDATION. OWASP Top Ten 2021. 2021. Disponível em: \url{https://owasp.org/www-project-top-ten/}.

  \item WAZLAWICK, R. S. \textbf{Metodologia de pesquisa para ciência da computação}. 2. ed. Rio de Janeiro: Elsevier, 2014.

  \item YAO, Shunyu et al. ReAct: Synergizing Reasoning and Acting in Language Models. In: \textit{Proceedings of the International Conference on Learning Representations (ICLR)}, 2023. Disponível em: \url{https://arxiv.org/abs/2210.03629}.

\end{itemize}

\end{document}
